% Pursuing Stacks (cotes 134-2 to 134-8): archive cote 134-7, site page 84 (PDF page 85), typescript page 526.
% typescript. Draft compared: ps6.tex:4220-4292.
% AI-transcribed by Claude Opus 5.5 (claude-opus-5-5, adaptive thinking, effort high), Message Batches API, batch ps4 (msgbatch_012Aoc1NuskATPsyBeKSYK7J), 2026-11-12; not yet checked by eye.
\scanpage{134-7/84}
\mspage{526}

Thus $\mathrm{L}\hand{\widetilde{\mathrm{H}}}_{\bullet}$ respects cofibration sequences, which means essentially
that for such a sequence in H\unsure{o}t$_{\circ}$(k)
\[ Y \hand{\longrightarrow} X \hand{\longrightarrow} Z \quad , \ Z \text{ the homotopy cofiber of } Y \hand{\rightarrow} X \ , \]
we get a corresponding exact triangle in Hotab$_{\circ}$(k)
\[ \text{(11 b)} \qquad
\diagram{\begin{tikzcd}[column sep=small]
 & \mathrm{L}\hand{\widetilde{\mathrm{H}}}_{\bullet}(Z) \ar[dl, "\hand{1}"'] & \\
 \mathrm{L}\hand{\widetilde{\mathrm{H}}}_{\bullet}(Y) \ar[rr] & & \mathrm{L}\hand{\widetilde{\mathrm{H}}}_{\bullet}(X) \ar[ul]
\end{tikzcd}} \qquad , \]
which is the most complete and elegant way, I guess, for expressing
behaviour of homology with respect to \xout{\unsure{susp}} homotopy cofibers and suspen-%
sions. DUally, $\mathrm{L}\hand{\Pi}_{\bullet}$ respects fibration sequences, which means that
for such a sequence in Hot$_{\circ}$(k)
\[ Z \hand{\rightarrow} X \hand{\rightarrow} Y \quad , \ Z \text{ the \inserted{(connected)} homotopy fiber of } X \hand{\rightarrow} Y \ , \]
we get a corresponding exact triangle in Hotab$_{\bullet}$(k)
\[ \text{(11 a)} \qquad
\diagram{\begin{tikzcd}[column sep=small]
 & \mathrm{L}\hand{\Pi}_{\bullet}(\unsure{\over{Z}{Y}}) \ar[dl, "\hand{1}"'] & \\
 \mathrm{L}\hand{\Pi}_{\bullet}(Z) \ar[rr] & & \mathrm{L}\hand{\Pi}_{\bullet}(X) \ar[ul]
\end{tikzcd}} \qquad , \]
which is a more complete and elegant way of stating the exact homotopy
sequence of a fibration, in terms of total homotopy.

\inserted[typed]{Formula} (11 b) implies that $\mathrm{L}\hand{\widetilde{\mathrm{H}}}_{\bullet}$ commutes with suspension functors $\hand{\Sigma}$ ,
the latter in Hotab$_{\bullet}$(k) (visualized as a derived chain complex category)
is just the shift functor
\[ \hand{\Sigma}_{ab} : \ L_{\bullet} \hand{\longmapsto} L_{\bullet}[1] \qquad ( L[1]_{n} = L_{n-1}) \quad \deleted{\unsure{$\ast$}} , \]
which is a \uline{fully faithful} functor . Dually, formula (11 a) implies
that $\mathrm{L}\hand{\Pi}_{\bullet}$ commutes to the loop-space functors $\hand{\Omega^{\circ}}$; the latter in
Hotab$_{\bullet}$(k) cannot be described as just a shift in opposite direction
\xout{\unsure{in}}
\[ L_{\bullet} \hand{\longmapsto} L_{\bullet}[-1] \ , \]
as this will get us out of \inserted[typed]{0-connected} chain complexes\unsure{\over{.}{,}} we have to truncate moreover
at dimension \unsure{1} (afterwards, or at dimension \xout{\illegible{1}} 2 beforehand). This func%
% NOTE: the right edge of the page is cut at the end of this line; a hyphen after "func" may be hidden
tor is not fully faithful therefore - we loose something when passing
from $L_{\bullet}$ to $\hand{\Omega^{\circ}}_{ab}(L_{\bullet})$.

As for the two functors $\hand{\widetilde{K}},\hand{\widetilde{S}}$ \xout{from} in the opposite direction, from
Hotab$_{\bullet}$(k) (i.e. essentially chain complexes) to Hot$_{\bullet}$(k), that the
first one respects fibration sequences is surely quite familiar a fact
(and kind of tautological) in the set-up of ordinary homotopy types. \inserted{That}
\over{T}{t}he second, less familiar functor $\hand{\widetilde{S}}$ should respect cofibration sequen-%

\margin{84}
